\documentclass[11pt,a4paper]{article}
% Shared preamble for the evaluation report.
\usepackage[T1]{fontenc}
\usepackage{lmodern}
\usepackage[english]{babel}
\usepackage[a4paper,margin=1in]{geometry}
\usepackage{booktabs}
\usepackage{graphicx}
\usepackage{array}
\usepackage{caption}
\usepackage{float}
\usepackage{amsmath}
\usepackage{amssymb}
\usepackage{longtable}
\usepackage{microtype}
\usepackage{enumitem}
\usepackage[table]{xcolor}
\usepackage{adjustbox}

% Shrink a table only if it would otherwise run into the margin. A table that
% fits keeps the document's type size; one that does not is scaled rather than
% left overflowing.
\newcommand{\fitwidth}[1]{\begin{adjustbox}{max width=\linewidth}#1\end{adjustbox}}
\usepackage[hidelinks]{hyperref}

\definecolor{ink}{HTML}{0B0B0B}
\definecolor{inksoft}{HTML}{52514E}
\definecolor{rule}{HTML}{DCDBD6}
\definecolor{ours}{HTML}{2A78D6}

\captionsetup{font=small,labelfont=bf,labelsep=period,skip=6pt}
\captionsetup[table]{position=top,skip=4pt}
\setlength{\tabcolsep}{5pt}
\renewcommand{\arraystretch}{1.12}

% A boxed statement of what a section establishes, set apart from the prose.
\newenvironment{finding}
  {\par\vspace{4pt}\noindent\begin{minipage}{\linewidth}%
   \color{ink}\leftskip=10pt\rightskip=10pt\small\itshape
   \vrule width 2pt height 0pt depth 0pt\hspace{-2pt}}
  {\par\end{minipage}\vspace{6pt}\par}

\newcommand{\ourwork}{\textcolor{ours}{\textbf{this work}}}


\newcommand{\resultfigures}{../../bench/results/figures/}
\newcommand{\resulttables}{../../bench/results/tables/}
\graphicspath{{../../bench/results/figures/}}
\input{../../bench/results/tables/facts}

\title{\bfseries An Experimental Evaluation of Persistence Strategies\\[2pt]
for Range-Add, Range-Sum Segment Trees\\[4pt]
{\large Exploratory pilot --- not a confirmatory study}}
\author{Version-Aware Lazy Segment Tree project}
\date{Generated from the exploratory pilot recorded in \texttt{bench/results/}}

\begin{document}
\maketitle

\noindent\fbox{\parbox{\dimexpr\linewidth-2\fboxsep-2\fboxrule}{\textbf{Status.}
This report describes an \emph{exploratory pilot}: one machine, one primary
compiler and allocator, run before any analysis protocol was registered. Its
purpose is to find harness defects, choose estimands and precision targets,
propose hypotheses and size a later confirmatory campaign. Nothing in it is a
confirmatory result, and none of its measurements will be pooled with, or
substituted for, the registered confirmatory campaign.}}
\medskip

% ===========================================================================
\section{Introduction}

A partially persistent segment tree keeps every version it has ever published
readable. This report measures \Fact{structureCount} ways of doing that for one
problem --- range additions on the newest version, range-sum queries against any
version --- and asks which is cheapest, and at what size the answer changes.

The design under study copies a path per update, as path copying always does, but
leaves the lazy tag on the copied node instead of pushing it into the children the
push would have to copy. Everything else is held fixed against a baseline
(\emph{copy-on-push}) that differs in that one decision alone, so the gap between
them is attributable to the decision rather than to the implementation.

\paragraph{What the campaign found.}
\begin{itemize}[leftmargin=1.4em,itemsep=1pt]
\item \textbf{Half the nodes.} \Fact{oursNodesPerUpdate} nodes stored per version
  against copy-on-push's \Fact{pushNodesPerUpdate} --- \Fact{pushNodeRatio}$\times$
  fewer, at \Fact{boundUsedShare}\% of a proved bound of \Fact{provedBound}. This is
  a count, so no machine can change it.
\item \textbf{Cheaper on both operations.} \Fact{updatecopyonpushMedian} on updates
  and \Fact{querycopyonpushMedian} on queries at the median cell; cheaper than every
  persistent rival in \Fact{rivalWins} of \Fact{rivalComparisons} comparisons, and in
  all \Fact{updatefatnodeCells} update comparisons against fat nodes and buffered
  slots.
\item \textbf{Flat in update width.} From one element to
  \Fact{widthPlateauTop} elements the update cost stays inside
  \Fact{widthPlateauSpread}$\times$; the widest update of all is the cheapest, at
  \Fact{widthWholeArray}\,$\mu$s.
\item \textbf{Cheapest at scale.} At one million elements it is the cheapest
  persistent structure on updates and on queries alike.
\item \textbf{Smallest memory.} \Fact{oursMemory}\,MiB where the closest competitor
  on time holds \Fact{memorycheckpointing}\,MiB.
\end{itemize}

Every persistent structure answered every query identically in every cell, so the
timing campaign doubles as a correctness campaign.

\begin{table}[H]
\centering
\caption{The structures. Bounds are per operation on $n$ elements after $U$
updates; $K$ is the checkpoint interval. The last keeps no history and is a floor,
not a competitor.}
\label{tab:structures}
\small
\begin{tabular}{lp{0.40\linewidth}ll}
\toprule
Name in the figures & Strategy & Update & Query \\
\midrule
This work & Path copying, lazy tag left in place on the copied node & $O(\log n)$ & $O(\log n)$ \\
Copy-on-push & Path copying, lazy tag pushed into freshly copied children & $O(\log n)$ & $O(\log n)$ \\
No lazy tags & Path copying to every element the update touches & $\Theta(k + \log n)$ & $O(\log n)$ \\
Full copy per version & A complete new tree for every version & $O(n)$ & $O(\log n)$ \\
Checkpoint and replay & Update log, full snapshot every $K$ versions & $O(\log n + n/K)$ & $O(\log n + K)$ \\
Buffered node slots & Two version-tagged slots per node, copy on overflow & $O(\log n)$ & $O(\log n)$ \\
Fat nodes & Several versioned states per node, copy on overflow & $O(\log n)$ & $O(\log n)$ \\
No history kept & Ordinary lazy segment tree, latest state only & $O(\log n)$ & $O(\log n)$ \\
\bottomrule
\end{tabular}
\end{table}

% ===========================================================================
\section{Setup}

\Fact{workloadCount} workloads, \Fact{structureCount} structures,
\Fact{totalCells} measurement cells, \Fact{totalTrials} recorded trials,
\Fact{totalOperations} operations. Each workload varies at most one property, so a
curve has one cause, and every structure replays a byte-identical stream.

\begin{table}[H]
\centering
\caption{The workloads.}
\label{tab:workloads}
\scriptsize
\fitwidth{\input{../../bench/results/tables/workloads}}
\end{table}

\begin{table}[H]
\centering
\caption{Everything a replication would have to match. The full procedure is in
\texttt{bench/README.md}.}
\label{tab:system}
\small
\fitwidth{\input{../../bench/results/tables/system}}
\end{table}

\paragraph{How the numbers are formed.}
The statistic is the median over \Fact{trialsPerCell} trials, each with its own
seed, with a percentile bootstrap interval from $10\,000$ resamples. Comparisons
are paired on the trial index, because every structure replays the same stream in
the same trial, and every reported ratio is the median of the per-trial ratios.
Significance is the Wilcoxon signed-rank test on the paired differences. The
\Fact{comparisonsTotal} comparisons form two correction families --- every
cell-by-baseline comparison for updates, and the same for queries --- and each
family is corrected separately by Benjamini--Hochberg at a 5\%
false-discovery rate.\footnote{Not Holm: the exact signed-rank test on eleven pairs
cannot report $p$ below $2^{-10} = \Fact{minWilcoxonP}$ however large the effect, so
family-wise control over several hundred comparisons would reject nothing at all.
Holm-adjusted values are computed and kept beside the reported ones in
\texttt{bench/results/summary/comparisons.csv}.}
Reported ratio bounds are the observed 2.5/97.5 percentiles of the per-trial
ratios, not confidence intervals. No outliers are removed anywhere.

% ===========================================================================
\section{Results}

\subsection{Scaling with array size}
\label{sec:scaling}

\begin{figure}[H]
\centering
\includegraphics[width=\linewidth]{scaling-update}
\caption{Time to publish one new version.}
\label{fig:scaling-update}
\end{figure}

\begin{figure}[H]
\centering
\includegraphics[width=\linewidth]{scaling-query}
\caption{Time to answer one historical range query, same runs.}
\label{fig:scaling-query}
\end{figure}

At \Fact{headlineSize} elements this work publishes a version in
\Fact{oursUpdate}\,$\mu$s and answers a historical query in \Fact{oursQuery}\,$\mu$s; at
\Fact{largestSize} elements, \Fact{oursUpdateBig}\,$\mu$s and \Fact{oursQueryBig}\,$\mu$s.
Publishing becomes \Fact{oursGrowth}$\times$ more expensive across the whole size sweep,
from \Fact{smallestSize} to \Fact{largestSize} elements. The tree's height only doubles over that span and the non-persistent
reference grows by \Fact{controlGrowth}$\times$, so most of the difference is the
memory hierarchy, not the algorithm: the arena outgrows each cache level in turn.

\begin{finding}
Every structure that copies nodes tracks the same shape and differs by a constant.
Checkpoint-and-replay does not: its curve bends upward and crosses this work
between one hundred thousand and one million elements on both operations.
\end{finding}

\begin{table}[H]
\centering
\caption{Every structure at \Fact{headlineSize} elements, balanced workload. CV is
the coefficient of variation across trials.}
\label{tab:headline}
\footnotesize
\fitwidth{\input{../../bench/results/tables/headline}}
\end{table}

\subsection{What the deferred tag buys}
\label{sec:tags}

\begin{figure}[H]
\centering
\includegraphics[width=\linewidth]{range-width}
\caption{Update cost against the number of elements one update covers.}
\label{fig:width}
\end{figure}

This experiment has a control at each end. At width one the tagless baseline does
exactly what this work does and costs \Fact{width1NoTags} as much --- marginally
less, carrying no tag arithmetic. As the width grows it must copy a path to every
element it touches: by \Fact{widthFullNoTagsAt} elements it costs
\Fact{widthFullNoTags} as much, and at \Fact{widthNoTagsRanOutAt} elements it stops
being measurable at all.

This work is flat over the same span --- \Fact{widthPlateauLow}\,$\mu$s to
\Fact{widthPlateauHigh}\,$\mu$s across four orders of magnitude --- because a fully
covered node absorbs the delta into a tag and its children are never visited. An
update covering the whole array is fully covered at the root, copies one node, and
costs \Fact{widthWholeArray}\,$\mu$s. \emph{The widest update in the sweep is by a wide
margin the cheapest.} The tag does not shrink the constant; it changes what the
update cost depends on.

Copy-on-push has the same shape here, as it must, since it also stops at a fully
covered node. The difference between the two policies is not visible on this axis.
It is visible in space.

\begin{figure}[H]
\centering
\includegraphics[width=\linewidth]{copies-per-update}
\caption{Nodes stored per published version --- a count, not a timing, so it does
not depend on the machine and can be checked against a proved bound.}
\label{fig:copies}
\end{figure}

\begin{finding}
Holding everything else identical, pushing tags into copied children costs
\Fact{pushNodeRatio}$\times$ the nodes per version and
\Fact{updatecopyonpushMedian} the update time at the median cell, over
\Fact{updatecopyonpushCells} comparable cells, of which this work is cheaper in
\Fact{updatecopyonpushWins}.
\end{finding}

\begin{table}[H]
\centering
\caption{Nodes stored per published version, by array size. The last row is the
worst case proved for this work, $4(h+1)$ for a tree of height $h$.}
\label{tab:space}
\small
\fitwidth{\input{../../bench/results/tables/space}}
\end{table}

\subsection{Depth of history, and shape of traffic}

\begin{figure}[H]
\centering
\includegraphics[width=\linewidth]{versions}
\caption{Cost against the number of versions retained, at a fixed array size.}
\label{fig:versions}
\end{figure}

This is the axis partial persistence exists for, and the one a structure has to be
flat on to be worth building. Path copying is flat because a version is a root and
reaching it costs nothing extra; the residual drift is the arena outgrowing cache.

\begin{figure}[H]
\centering
\includegraphics[width=\linewidth]{traffic-sweeps}
\caption{Four properties of the traffic, each varied on its own axis.}
\label{fig:sweeps}
\end{figure}

Reads concentrated on recent versions get cheaper for everyone, since those
versions are still in cache. Updates that change nothing are answered by
republishing the existing root, so cost falls as their share rises. Updates
concentrated in one region get cheaper for this work --- the copied paths overlap
and stay resident --- while copy-on-push climbs steeply, because a hot region is
exactly where tags accumulate and every accumulated tag is a push it pays for.

The fourth panel was included expecting it to be hostile to checkpoint-and-replay,
and it is not; reporting that is the point of including it. Replay always starts at
the nearest snapshot at or before the requested version, so age costs it nothing.
What the panel does show is that concentrating reads anywhere makes them cheaper
for every structure. Confined to the oldest \Fact{auditOldestShare}\% of history
this work answers in \Fact{auditOldestOurs}\,$\mu$s and checkpoint-and-replay costs
\Fact{auditOldestcheckpointing} that.

\subsection{Memory, and what it makes impossible}
\label{sec:feasibility}

\begin{figure}[H]
\centering
\includegraphics[width=\linewidth]{scaling-memory}
\caption{Memory retained after two hundred thousand operations.}
\label{fig:memory}
\end{figure}

\begin{figure}[H]
\centering
\includegraphics[width=\linewidth]{feasibility}
\caption{How much of the campaign each structure could run inside the memory limit.}
\label{fig:feasibility}
\end{figure}

A ceiling is a result, not a missing row. Full copy per version finished
\Fact{finishedfullcopy} of \Fact{workloadsTotal} workloads and the tagless baseline
\Fact{finishedpointonly}. Checkpoint-and-replay finished
\Fact{finishedcheckpointing}: its snapshots grow as $nU/K$ rather than with what
could have been shared, so at \Fact{headlineSize} elements it already holds
\Fact{memorycheckpointing}\,MiB against this work's \Fact{oursMemory}\,MiB. This
work, copy-on-push, buffered slots and fat nodes finished all twelve.

\begin{figure}[H]
\centering
\includegraphics[width=\linewidth]{space-time}
\caption{Speed against memory. Bottom left is better on both.}
\label{fig:pareto}
\end{figure}

\subsection{The slowest operations, not the typical one}

\begin{figure}[H]
\centering
\includegraphics[width=\linewidth]{latency-tail}
\caption{The distribution of a single update, median to slowest observed.}
\label{fig:tail}
\end{figure}

At \Fact{headlineSize} elements the median single update costs
\Fact{oursTailP50}\,$\mu$s and the slowest one in a hundred \Fact{oursTailP99}\,$\mu$s. The
extreme values --- up to \Fact{oursTailMax}\,$\mu$s --- are the \texttt{std::vector}
arena doubling and copying itself, not the persistence strategy; a structure that
retains more nodes pays that more often and at greater size.

\subsection{Head to head}
\label{sec:headtohead}

\begin{table}[H]
\centering
\caption{Every comparison in the campaign, one row per baseline. ``Cells'' counts
those where both structures completed every paired trial; ``separated'' counts
those surviving Benjamini--Hochberg correction at a 5\% false-discovery rate;
``unanimous'' counts those where every paired trial agreed.}
\label{tab:headtohead}
\small
\fitwidth{\input{../../bench/results/tables/head-to-head}}
\end{table}

Across \Fact{rivalComparisons} comparisons against another persistent structure,
this work is cheaper in \Fact{rivalWins}. It wins every one of the
\Fact{updatefatnodeCells} update comparisons against fat nodes and against buffered
node slots, by a median of \Fact{updatefatnodeMedian} and
\Fact{updatebufferedMedian}, and never loses either.

\begin{figure}[H]
\centering
\includegraphics[width=\linewidth]{speedup-matrix}
\caption{The whole comparison, cell by cell.}
\label{fig:matrix}
\end{figure}

% ===========================================================================
\section{Against checkpoint-and-replay: which to reach for}
\label{sec:checkpoint}

Checkpoint-and-replay is the only baseline that beats this work anywhere it
matters, so the boundary between them is worth drawing precisely. It reads the
nearest snapshot at or before the requested version and, for each logged update
between them, adds the overlap of that update's range with the queried range times
its delta --- a few arithmetic operations on contiguous memory rather than a tree
descent. At small $n$ that is genuinely hard to beat.

\begin{figure}[H]
\centering
\includegraphics[width=\linewidth]{checkpoint-tradeoff}
\caption{Checkpoint-and-replay swept across its checkpoint interval, against this
work, which has no such parameter. It is run at $K = \sqrt{n}$ everywhere else in
this report --- the value balancing its $O(\log n + n/K)$ update against its
$O(\log n + K)$ query, chosen from the bound rather than by hand.}
\label{fig:checkpoint}
\end{figure}

\begin{table}[H]
\centering
\caption{Where each strategy wins.}
\label{tab:versus}
\small
\begin{tabular}{p{0.30\linewidth}p{0.29\linewidth}p{0.29\linewidth}}
\toprule
Condition & Checkpoint and replay & This work \\
\midrule
Array size & cheaper up to $\approx10^5$ elements & cheaper at
\Fact{largestSize} elements, by \Fact{marginAtLargestUpdate}$\times$ on updates and
\Fact{marginAtLargestQuery}$\times$ on queries \\
Memory & \Fact{memorycheckpointing}\,MiB at \Fact{headlineSize} elements &
\Fact{oursMemory}\,MiB --- \Fact{memoryRatiocheckpointing}$\times$ less \\
Reaching the whole campaign & finished \Fact{finishedcheckpointing} of
\Fact{workloadsTotal} workloads & finished all \Fact{workloadsTotal} \\
Tuning & one knob; outside a narrow band of $K$ one operation collapses
(Figure~\ref{fig:checkpoint}) & no equivalent knob \\
Other aggregates & advantage vanishes: projecting an update onto a query range
needs both to be linear & indifferent to which aggregate is stored \\
Handles to past versions & cannot give one; it reconstructs on demand & a version
is a root, shareable and walkable \\
\bottomrule
\end{tabular}
\end{table}

\begin{figure}[H]
\centering
\includegraphics[width=\linewidth]{versus-checkpoint}
\caption{Every comparison against checkpoint-and-replay, grouped by the condition
that decides it. Each condition is a property of the workload definition, so no
cell was placed by looking at its result.}
\label{fig:versus}
\end{figure}

Figure~\ref{fig:versus} shows that the two do not trade cells at random. This work
takes every one of the five cells where an update covers the whole array --- by up
to $26.6\times$ on updates and $18.2\times$ on queries, because such an update copies
one node here and forces a full snapshot there. It takes every update cell at one
million elements, and every update cell once a hundred thousand versions have
accumulated, by a median of $3.15\times$. On queries it takes every cell of
query-heavy traffic and every cell of skewed or historical reads --- the audit case
--- because checkpointing pays its replay again on each query while a version here
is just a root.

Checkpointing's advantage lives in one place. It wins
\Fact{lossesWorstCount} cells in the whole campaign, and 39 of them are small
arrays with narrow updates --- 39 of the 42 comparisons in that group --- where
its projection is a handful of cache-resident additions. It needs all three
conditions at once: a small array, a linear aggregate, and a well-chosen $K$.
Break any one --- let the array grow, make the aggregate a minimum, or ask for a
handle to an old version --- and this work is the better choice.

% ===========================================================================
\section{Trade-offs and limitations}

\paragraph{What the design costs.}
\begin{itemize}[leftmargin=1.4em,itemsep=1pt]
\item \emph{The win is a constant factor.} Against copy-on-push both are
  $O(\log n)$ per operation. The result is a factor of two in space and
  \Fact{updatecopyonpushMedian} in time, not a better bound.
\item \emph{History is not free.} Against a tree that retains nothing this work
  pays a median of \Fact{persistenceCostMedian}$\times$ per update, over a range of
  \Fact{persistenceCostRange}$\times$.
\item \emph{Fat nodes and buffered slots store fewer nodes per version}
  (Table~\ref{tab:space}); this work wins on time and on total memory, not on node
  count.
\item \emph{The tag is overhead when nothing is deferred.} At single-element
  updates the tagless baseline is \Fact{width1NoTags} the cost.
\item \emph{Documented payload is not what is paid.} A growing
  \texttt{std::vector} holds up to twice its elements' bytes. At
  \Fact{headlineSize} elements this work requests \Fact{allocOverheadOurs}$\times$
  its documented payload at peak, and no persistent structure exceeds
  \Fact{allocOverheadMax}$\times$. The space comparisons here use the documented
  payload, which is the conservative choice: this work stores the fewest nodes and
  so has the most to gain from the other measure.
\item \emph{Scope.} Range add and range sum, updates to the newest version only.
  Branching updates, full persistence and non-linear aggregates are out of scope.
\end{itemize}

\begin{table}[H]
\centering
\caption{Checks on the measurement itself.}
\label{tab:validity}
\small
\fitwidth{\input{../../bench/results/tables/validity}}
\end{table}

Every number here comes from one \Fact{machine} with one compiler, and every
stream is synthetic. The node counts in Table~\ref{tab:space} and the asymptotic
separations do not depend on the machine; the close constant-factor results do.

The machine has a warm-up transient large enough to matter --- on an early
configuration around $1.7\times$ on the first recorded trials, confined to updates
and scaling with how much a structure allocates. It decays with elapsed load rather
than with replay count, so the harness runs a fifteen-second duration-based warm-up
before recording anything and discards three further trials per cell. That leaves a
residual of \Fact{warmupResidual}$\times$ on the pooled first trial, this work
varying by \Fact{oursMedianCV}\% across its trials, and no order effect on the
medians (\Fact{orderEffect}\% across all eight positions). Pairing on the trial
index cancels whatever remains.

% ===========================================================================
\section{Conclusion}

Across \Fact{totalCells} cells and \Fact{totalTrials} trials, leaving a lazy tag on
the copied node instead of pushing it into copied children is cheaper in both time
and space, and the space result --- \Fact{pushNodeRatio}$\times$ fewer nodes per
version --- is a count that no machine can change. Against the two in-node
versioning strategies this work is cheaper on every update comparison in the
campaign. Against strategies that do not share structure it is cheaper by
\Fact{updatepointonlyMedian} and \Fact{updatefullcopyMedian} at the median and, more
to the point, finishes problems they cannot start.

Against checkpoint-and-replay it loses at small sizes and wins at large ones, and
Table~\ref{tab:versus} says which to reach for. Path copying with a deferred tag is
the choice whenever the array is large, the aggregate is not linear, or past
versions must exist as objects rather than as answers.

\end{document}
